\documentclass[11pt]{article}
\usepackage[a4paper,margin=24mm]{geometry}
\usepackage{amsmath,amssymb,booktabs}
\setlength{\parskip}{0.4em}\setlength{\parindent}{0pt}
\newcommand{\Wz}{W\!z_{CoM}}
\newcommand{\te}{\tau_{\rm end}}
\newcommand{\rmax}{\rho_{\max}}
\newcommand{\rbar}{\bar\rho}
\title{\vspace{-18mm}\textbf{How (95)--(97) are obtained}\\[2mm]
\large a step-by-step companion to Sec.~VI-E}
\date{\small every step below is checked by
\texttt{analysis/vi\_e\_geometric\_box.py}}
\begin{document}\maketitle

\section*{0. What is being bounded, and why the rate}

The estimator does not see the tilt angle. It fits the two-segment
model to the \emph{rate}
\begin{equation}
  \omega_{\rm nom}(\tau) = C_1\bigl(\cosh C_2\tau - 1\bigr),
  \qquad C_1 = \frac{\dot M}{\Wz},
  \qquad C_2 = \sqrt{\frac{\Wz}{J_P}},
  \tag{D.1}
\end{equation}
reconstructed from the gyro, and reports the onset as the argument
that minimises the residual of that fit. Two questions follow, and
(97) answers both:

\begin{enumerate}
\item[(i)] how far can the \emph{true} rate depart from (D.1) before
  the fit stops describing the data? --- the quantity
  $|e_\omega|_{\rm end}/\omega_{\rm nom,end}$;
\item[(ii)] how far does that departure move the \emph{reported
  onset moment}? --- the quantity $|\Delta M_{\rm crit}|$.
\end{enumerate}

These are different questions with different answers, which is why the
two halves of (97) carry different constants ($\tfrac12$ and $1$ in the
first, $\tfrac17$ and $\tfrac15$ in the second). Section 5 answers (i)
and Section 6 answers (ii). If only one line of this document is read,
let it be that one.

The departure obeys the Duhamel integral (88),
\begin{equation}
  e_\omega(\tau) = \int_0^{\tau}\dot h_\varphi(\tau-s)\,\rho(s)\,ds ,
  \qquad
  \dot h_\varphi(u) = \frac{\cosh C_2 u}{J_P},
  \tag{D.2}
\end{equation}
with $\rho$ the part of the true forcing that the estimator's span
$\{1,\tau,\delta\varphi\}$ cannot absorb --- the gravity remainder
$\rho_\varphi$ and the bilinear ground-effect coupling $\rho_{GE}$ of
(89). Everything below is about turning (D.2) into a number.

\section*{1. The crude bound, and why it is not enough}

Put $|\rho(s)|\le\rmax$ inside (D.2) and pull it out:
\begin{equation}
  |e_\omega(\te)| \;\le\; \rmax\!\int_0^{\te}\!\dot h_\varphi(\te-s)\,ds
  \;=\; \rmax\!\int_0^{\te}\!\frac{\cosh C_2u}{J_P}\,du
  \;=\; \frac{\rmax\sinh x}{J_PC_2},
  \quad x \triangleq C_2\te .
  \tag{D.3}
\end{equation}
The substitution $u = \te - s$ was used in the middle step; the
integral of $\cosh$ is $\sinh$. Using $J_P = \Wz/C_2^{2}$, i.e.\
$J_PC_2 = \Wz/C_2$, the inertia disappears:
$|e_\omega(\te)| \le \rmax C_2\sinh x/\Wz$. This is (90).

(D.3) is honest but blunt. It treats $\rho$ as though it sat at its
maximum for the whole window, whereas both channels start at zero at
the onset and grow like $\tau^6$ and $\tau^4$. Evaluated over the
geometric box of Sec.~VI-D it returns a deviation of $446\%$ at the
slowest ramp --- bigger than the signal it is supposed to bound. Sections 2--4
recover the lost factor.

\section*{2. Chebyshev's integral inequality}

\textbf{Statement.} Let $f,g$ be integrable on $[a,b]$. If they are
\emph{oppositely} monotone (one non-decreasing, the other
non-increasing), then
\begin{equation}
  \int_a^b\! fg \;\le\; \frac{1}{b-a}\int_a^b\! f \int_a^b\! g .
  \tag{D.4}
\end{equation}
If they are monotone in the \emph{same} direction, the inequality
reverses. (Textbooks usually print the same-direction version, with
$\ge$; the one needed here is the reversed one, so the direction is
worth checking rather than remembering.)

\textbf{Proof.} Opposite monotonicity means
$\bigl(f(u)-f(v)\bigr)\bigl(g(u)-g(v)\bigr)\le0$ for every
$u,v\in[a,b]$: when $u>v$ one bracket is $\ge0$ and the other $\le0$.
Integrate that product over the square $[a,b]^2$ and expand the four
terms,
\[
  \underbrace{\int\!\!\int f(u)g(u)}_{(b-a)\int fg}
  - \underbrace{\int\!\!\int f(u)g(v)}_{\int f\int g}
  - \underbrace{\int\!\!\int f(v)g(u)}_{\int f\int g}
  + \underbrace{\int\!\!\int f(v)g(v)}_{(b-a)\int fg}
  \;=\; 2(b-a)\!\int\! fg - 2\!\int\! f\!\int\! g \;\le\; 0,
\]
which is (D.4). $\square$

\textbf{Application to (D.2).} Take $[a,b] = [0,\te]$ with
$f(s) = \rho(s)$ and $g(s) = \dot h_\varphi(\te-s)$.

\begin{itemize}
\item $f$ is non-decreasing: both channels vanish at the onset and
  grow monotonically over the window ($\rho_\varphi\propto\varphi^2$
  with $\varphi$ rising, $\rho_{GE}\propto\tau\varphi$ likewise).
  Equivalently, both are power series in $u$ with non-negative
  coefficients, so $f'\ge0$ term by term --- the same sign pattern that
  Sec.~3 uses for a different purpose. This is the one place the
  argument needs the channels to increase; the reduction factors of
  Sec.~3 do not.
\item $g$ is non-increasing \emph{in $s$}: as $s$ grows the argument
  $\te-s$ shrinks and $\cosh$ decreases. Note the kernel is increasing
  in its own argument; it is the reflection $s\mapsto\te-s$ that makes
  it decreasing in $s$.
\end{itemize}

They are oppositely monotone, so (D.4) applies with $\le$, which is the
useful direction:
\begin{equation}
  e_\omega(\te)
  \;\le\; \underbrace{\frac{1}{\te}\int_0^{\te}\!\rho}_{\textstyle\rbar}
  \;\cdot\;\int_0^{\te}\!\dot h_\varphi(\te-s)\,ds
  \;=\; \frac{\rbar\sinh x}{J_PC_2}.
  \tag{D.5}
\end{equation}
(D.5) is (D.3) with $\rmax$ replaced by the \emph{time average} $\rbar$.
Nothing else changed. All that remains is to compute $\rbar/\rmax$ for
the two channels.

\section*{3. The reduction factors $R_\varphi$ and $R_{GE}$: (95)}

Both channels have a known shape, because $\varphi$ does. From (D.1),
$\varphi_{\rm nom}(\tau) = (C_1/C_2)\Lambda(C_2\tau)$ with
$\Lambda(u) = \sinh u - u$. Hence, writing $u = C_2\tau$,
\[
  \rho_\varphi(\tau) = \tfrac12 G''\varphi^2 = A\,\Lambda(u)^2,
  \qquad
  \rho_{GE}(\tau) = \beta_M\dot M\tau\,\varphi = B\,u\,\Lambda(u),
\]
with $A,B$ constants whose values never matter --- they cancel in the
ratio below.

\textbf{Gravity channel.} The maximum is at the window end,
$\rho_{\varphi,\max} = A\Lambda(x)^2$. The average is, substituting
$u = C_2\tau$ so that $d\tau = du/C_2$ and $\te = x/C_2$,
\[
  \bar\rho_\varphi = \frac{1}{\te}\int_0^{\te}\! A\Lambda(C_2\tau)^2 d\tau
  = \frac{C_2}{x}\cdot\frac{A}{C_2}\int_0^x\!\Lambda(u)^2 du
  = \frac{A}{x}\int_0^x\!\Lambda(u)^2 du ,
\]
so the constant $A$ cancels and
\begin{equation}
  R_\varphi(x) \;\triangleq\;
  \frac{\bar\rho_\varphi}{\rho_{\varphi,\max}}
  \;=\; \frac{1}{x}\,\frac{\int_0^x\Lambda(u)^2du}{\Lambda(x)^2}.
  \tag{D.6}
\end{equation}

\textbf{Ground-effect channel.} Identically,
$\rho_{GE,\max} = B\,x\Lambda(x)$ and
\begin{equation}
  R_{GE}(x) \;=\; \frac{1}{x}\,
  \frac{\int_0^x u\,\Lambda(u)\,du}{x\,\Lambda(x)} .
  \tag{D.7}
\end{equation}

\textbf{How to read $1/7$ and $1/5$.} The quickest way to see where
these come from --- not the proof, which is below --- is the order of
each channel. Keeping the leading term of
$\Lambda(u) = u^3/6 + u^5/120 + \cdots$ makes $\Lambda^2\simeq u^6/36$
and $u\Lambda\simeq u^4/6$, so the two channels are sixth and fourth
order near the onset, and a $\tau^n$ signal averages over $[0,T]$ to
$T^n/(n{+}1)$ against its endpoint value $T^n$: hence $\tfrac17$ and
$\tfrac15$. At the other end $\Lambda\to\sinh u\simeq e^u/2$ gives
$\int_0^x\Lambda^2\simeq e^{2x}/8$ and $\int_0^xu\Lambda\simeq xe^x/2$,
so $R_\varphi\simeq1/(2x)$ and $R_{GE}\simeq1/x$ --- the origin of the
$\tfrac12$ and $1$ that Sec.~5 needs, once $\Psi\simeq x$ multiplies
them.

Neither reading is used as a proof. (D.6b) below obtains the small-$x$
limits from l'H\^opital's rule with no expansion of $f$ whatever, and
Sec.~5 obtains the large-$x$ ones the same way. Only the monotonicity
argument that follows uses a series, and only for its sign pattern.

\textbf{Why the small-$x$ value is the supremum.} Both $R$'s decrease
throughout, and this needs no numerics. Substituting $u = xt$ in the
numerator rescales the window to $[0,1]$:
\begin{equation}
  R(x) = \frac{1}{x}\frac{\int_0^x f(u)\,du}{f(x)}
       = \int_0^1 \frac{f(xt)}{f(x)}\,dt .
  \tag{D.6a}
\end{equation}
Let
\[
  n_{\rm eff}(u) \;\triangleq\; \frac{u f'(u)}{f(u)}
  \;=\; \frac{d\log f}{d\log u}
\]
be the \emph{local order} of $f$. The small-$x$ limit of $R$ follows
from it by one application of l'H\^opital's rule, with no expansion of
$f$: in the first form of (D.6a) numerator and denominator both vanish
at $x=0$, so
\begin{equation}
  \lim_{x\to0}R(x)
  = \lim_{x\to0}\frac{\int_0^x\!f(u)\,du}{x\,f(x)}
  = \lim_{x\to0}\frac{f(x)}{f(x) + x f'(x)}
  = \frac{1}{1 + n_{\rm eff}(0^+)} ,
  \tag{D.6b}
\end{equation}
the last step dividing through by $f$. Only the limit of the local
order is needed, never the coefficients of $f$. Differentiating the
integrand of
(D.6a) logarithmically,
\[
  \frac{\partial}{\partial x}\log\frac{f(xt)}{f(x)}
  = t\,\frac{f'(xt)}{f(xt)} - \frac{f'(x)}{f(x)}
  = \frac{n_{\rm eff}(xt) - n_{\rm eff}(x)}{x} .
\]
For $t\in(0,1)$ we have $xt < x$, so if $n_{\rm eff}$ is increasing the
bracket is $\le0$: every integrand of (D.6a) decreases in $x$, hence so
does $R$. $\square$

Both channels satisfy the hypothesis, and with
$\Lambda' = \cosh u - 1$ the local orders are expressible through the
same $\Psi_\varphi(u) = u(\cosh u - 1)/(\sinh u - u) = u\Lambda'/\Lambda$
that Sec.~4 discarded as the wrong normaliser:
\[
  f = \Lambda^2:\quad n_{\rm eff} = 2\Psi_\varphi(u)
  \;\; \text{from } 6 \text{ to } \infty;
  \qquad
  f = u\Lambda:\quad n_{\rm eff} = 1 + \Psi_\varphi(u)
  \;\; \text{from } 4 \text{ to } \infty.
\]

\textbf{Why $n_{\rm eff}$ increases.} What is required here is easy to
understate, so it is worth naming precisely. Since $n_{\rm eff} =
d\log f/d\log u$, the claim
\[
  \frac{dn_{\rm eff}}{d\log u}
  = \frac{d^2(\log f)}{d(\log u)^2} \;\ge\; 0
\]
says that $\log f$ is \emph{convex} in $\log u$ --- that the slope of
$f$ on log--log axes grows. Monotonicity of $f$ gives only
$n_{\rm eff}\ge0$, which is the slope's \emph{sign} and nothing more,
and the two properties are independent: $f(u) = 1-e^{-u}$,
$u/(1{+}u)$ and $\log(1{+}u)$ are all strictly increasing and all have
$n_{\rm eff}$ strictly \emph{decreasing}, hence $R$ increasing. A
saturating signal has a higher average relative to its endpoint the
longer the window, not a lower one. The two channels here behave the
opposite way not because they rise but because they \emph{stiffen}.

Convexity does hold, and one place a power series is used --- only for
its sign pattern, with no coefficient evaluated --- establishes it.

\emph{Lemma (log--log convexity).} Let $f(u) = \sum_{n\ge0}a_nu^n$ converge for all $u\ge0$
with every $a_n\ge0$, at least two of them non-zero. Then
$n_{\rm eff} = uf'/f$ is strictly increasing on $(0,\infty)$.

The second condition is what makes $f>0$ for $u>0$, so that dividing by
$f$ below is legitimate, and what makes the conclusion strict. No
monotonicity of $f$ is assumed --- it would be redundant, since
$f' = \sum_{n\ge1}na_nu^{n-1}\ge0$ term by term, and it is not used in
the proof.

\emph{Proof.} Differentiating $n_{\rm eff}$,
\[
  u\,\frac{dn_{\rm eff}}{du}
  = \frac{uf'}{f} + \frac{u^2f''}{f}
    - \Bigl(\frac{uf'}{f}\Bigr)^{\!2}
  = \frac{S_0S_2 - S_1^2}{S_0^{\,2}},
  \qquad
  S_k \triangleq \sum_n n^k a_nu^n ,
\]
since $uf' = S_1$ and $u^2f'' = S_2 - S_1$, so the first two terms
combine to $S_2/S_0$. Every $S_k$ runs from $n = 0$, including the two
that come from differentiating: $f'$ starts at $n=1$ and $f''$ at
$n=2$, but the coefficients $n$ and $n(n-1) = n^2 - n$ vanish on
precisely the terms those shifts drop, so extending all three sums to
the same index set neither adds nor loses anything. It is worth doing,
because Cauchy--Schwarz is about to treat them as sequences over one
index set, and those sequences are infinite: $\Lambda$, and with it
$\Lambda^2$ and $u\Lambda$, has infinitely many non-zero coefficients.
That is what the hypothesis of convergence for every $u\ge0$ is for.
A power series converging at some $u_0>0$ converges absolutely for
$|u|<u_0$, so convergence at every $u\ge0$ makes the radius infinite,
and the termwise derivatives $\sum na_nu^{n-1}$ and
$\sum n(n{-}1)a_nu^{n-2}$ then converge on the same disc. Hence $f''$
exists as written and $S_1$, $S_2$ are finite --- the two sequences lie
in $\ell^2$, which is what Cauchy--Schwarz for infinite sums requires.

Cauchy--Schwarz applied to $\bigl\{\sqrt{a_nu^n}\bigr\}$ and
$\bigl\{n\sqrt{a_nu^n}\bigr\}$ --- the square roots real because
$a_nu^n\ge0$, which is the only place the sign pattern is needed ---
gives $S_1^2\le S_0S_2$. Hence
$dn_{\rm eff}/du\ge0$, with equality only when the two sequences are
proportional, i.e.\ only when a single $a_n$ is non-zero. $\square$

(Dividing through by $S_0$ turns $S_k/S_0$ into the moments of the
weights $a_nu^n/f$ and the numerator into their variance, which is a
shorter way to say the same thing; nothing below depends on reading it
that way.)

Both kernels qualify: $\Lambda = \sum_{k\ge1}u^{2k+1}/(2k+1)!$ has
non-negative coefficients, and so do the product $\Lambda^2$ and the
shift $u\Lambda$; each has infinitely many non-zero, so certainly two.
Hence $n_{\rm eff}$ is strictly increasing for both channels and $R$ is
strictly decreasing. Reading the $a_nu^n/f$ as weights also returns the
origin values, since as $u\to0$ they concentrate on the lowest exponent
present, $6$ and $4$ --- the same numbers the l'H\^opital route below
obtains without any series.

\textbf{The same fact as a closed-form derivative.} An appendix that
quotes $dR/dx$ explicitly needs those expressions signed, and the
argument above does it with no further machinery --- in particular
without the twelve-stage differentiation that signing them directly
would cost. Write $A(x) = \int_0^xf$, so $R = A/(xf)$. Then, using
$A' = f$ and $xf' = n_{\rm eff}f$,
\begin{equation}
  R'(x) = \frac{A'\,xf - A\,(f + xf')}{x^2f^2}
        = \frac{x f(x) - \bigl(1 + n_{\rm eff}(x)\bigr)A(x)}{x^2 f(x)} .
  \tag{D.7a}
\end{equation}
The sign turns on whether $A$ exceeds $xf/(1+n_{\rm eff})$, and it does,
for exactly the reason $n_{\rm eff}$ increases. Since
$d\log f/d\log u = n_{\rm eff}$ and $n_{\rm eff}(t)\le n_{\rm eff}(x)$
for $t\le x$,
\[
  \log\frac{f(x)}{f(u)} = \int_u^x n_{\rm eff}(t)\,\frac{dt}{t}
  \;\le\; n_{\rm eff}(x)\,\log\frac{x}{u}
  \quad\Longrightarrow\quad
  f(u) \;\ge\; f(x)\Bigl(\frac{u}{x}\Bigr)^{n_{\rm eff}(x)} ,
\]
and integrating that minorant over $[0,x]$,
\[
  A(x) \;\ge\; f(x)\!\int_0^x\!\Bigl(\frac{u}{x}\Bigr)^{n_{\rm eff}(x)}\!du
       \;=\; \frac{x\,f(x)}{1 + n_{\rm eff}(x)} ,
\]
strictly, since $n_{\rm eff}$ is not constant. The numerator of (D.7a)
is therefore negative and $R'<0$. $\square$

In words: $f$ is at least as steep as the power law of its local order
at $x$, so its average over the window is at least the $1/(n{+}1)$ that
power law would give --- and $1/(n{+}1)$ is precisely the value $R$
would need for $R'$ to vanish.

Substituting $n_{\rm eff} = 2x\Lambda'/\Lambda$ for $f = \Lambda^2$, and
$n_{\rm eff} = (\Lambda + x\Lambda')/\Lambda$ for $f = u\Lambda$, and
clearing denominators puts (D.7a) in the form an appendix would quote:
\begin{equation}
  \frac{dR_\varphi}{dx}
  = \frac{x\Lambda^3 - (\Lambda + 2x\Lambda')\!\int_0^x\!\Lambda^2du}
         {x^2\Lambda^3},
  \quad
  \frac{dR_{GE}}{dx}
  = \frac{x^2\Lambda^2 - (2\Lambda + x\Lambda')\!\int_0^x\!u\Lambda\,du}
         {x^3\Lambda^2},
  \tag{D.7b}
\end{equation}
both negative for $x>0$. The factors of $x$ inside the brackets are
load-bearing: with them the numerators cancel to order $x^{12}$ and
$x^{10}$ at the origin, which is what leaves $R'$ finite there; drop
either and the expression diverges as $x^{-2}$ instead. Both forms are
checked against a five-point difference of $R$ by
\texttt{analysis/closed\_form\_check.py}.

\textbf{The origin values without the series.} The route above needs
the coefficients' sign pattern; this one needs nothing, and is worth
recording because it settles the same two numbers independently. The
value of $\Psi_\varphi$ at the origin is l'H\^opital's rule, applied
three times. Writing $\Psi_\varphi = N/D$
with $N = u\cosh u - u$ and $D = \sinh u - u$, successive
differentiation gives the four pairs
\[
  \frac{u\cosh u - u}{\sinh u - u},\quad
  \frac{\cosh u + u\sinh u - 1}{\cosh u - 1},\quad
  \frac{2\sinh u + u\cosh u}{\sinh u},\quad
  \frac{3\cosh u + u\sinh u}{\cosh u},
\]
of which the first three are $0/0$ at $u = 0$ while the fourth is
$3/1$. Hence $\Psi_\varphi(0^+) = 3$, so $n_{\rm eff}(0^+)$ is $6$ for
$\Lambda^2$ and $4$ for $u\Lambda$, and (D.6b) returns $\tfrac17$ and
$\tfrac15$. The mechanism is physical: $\Lambda$ behaves as $u^3$ near the onset and
grows exponentially later, so the larger the window the more sharply
both channels are concentrated at its end, and the smaller their
average is relative to the endpoint value. The small-$x$ values are
therefore suprema:
\begin{equation}
  \rbar = R_\varphi(x)\,\rho_{\varphi,\max}
        + R_{GE}(x)\,\rho_{GE,\max},
  \qquad R_\varphi\le\tfrac17,\qquad R_{GE}\le\tfrac15 .
  \tag{95}
\end{equation}
The suprema are approached but never attained: for $x>0$ the
inequalities are strict. Over the flown $x\in[1.79,5.20]$ the factors
run $0.089$--$0.133$ and $0.145$--$0.191$, so quoting $\tfrac17$ and
$\tfrac15$ gives away $10$--$60\%$ --- while gaining a factor of $5$ to
$11$ over (D.3).

\section*{4. Normalising on the rate: (96)}

Divide (D.5) by the nominal rate at the same instant,
$\omega_{\rm nom,end} = C_1(\cosh x - 1)$ with $C_1 = \dot M/\Wz$, and
use $J_PC_2 = \Wz/C_2$:
\[
  \frac{|e_\omega|_{\rm end}}{\omega_{\rm nom,end}}
  \;\le\; \frac{\rbar\,C_2\sinh x/\Wz}
               {(\dot M/\Wz)(\cosh x - 1)}
  \;=\; \frac{\rbar\,C_2}{\dot M}\cdot\frac{\sinh x}{\cosh x - 1}.
\]
$\Wz$ has cancelled. Now eliminate $\dot M$ in favour of the moment
increment the window actually spans. Since $\te = x/C_2$,
\[
  \Delta M_{\rm win} = \dot M\te = \frac{\dot M x}{C_2}
  \qquad\Longleftrightarrow\qquad
  \frac{C_2}{\dot M} = \frac{x}{\Delta M_{\rm win}},
\]
and substituting,
\begin{equation}
  \frac{|e_\omega|_{\rm end}}{\omega_{\rm nom,end}}
  \le \frac{\Psi(x)\,\rbar}{\Delta M_{\rm win}},
  \qquad
  \Psi(x) = \frac{x\sinh x}{\cosh x - 1} .
  \tag{96}
\end{equation}
The half-angle identities $\cosh x - 1 = 2\sinh^2(x/2)$ and
$\sinh x = 2\sinh(x/2)\cosh(x/2)$ collapse the quotient:
\[
  \Psi(x) = x\,\frac{2\sinh(x/2)\cosh(x/2)}{2\sinh^2(x/2)}
          = x\coth\frac{x}{2} \;\ge\; 2 .
\]
$\Psi$ rises from $2$ at $x\to0$ to $x$ as $x\to\infty$. (Had the
excursion been used as the normaliser instead of the rate, the factor
would have been $x(\cosh x - 1)/(\sinh x - x)$, rising from $3$ ---
a different function, and the wrong one, since the estimator never
looks at $\varphi$.)

\section*{5. Why $\Psi\rbar$ stays bounded: first half of (97)}

$\Psi$ grows with $x$ and the $R$'s shrink with it, and the two
products converge to $\tfrac12$ and $1$. Both bounds are exact, and
neither needs the series: one corollary of the mean value theorem
settles both, three applications for the first and one for the
second.

\textbf{The claims in integral form.} Since
$\Psi = x\coth(x/2)$ and $\rho_{\varphi,\max} = A\Lambda(x)^2$,
\[
  \Psi R_\varphi
  = x\coth\tfrac{x}{2}\cdot\frac{\int_0^x\Lambda^2du}{x\Lambda(x)^2}
  = \coth\tfrac{x}{2}\,\frac{\int_0^x\Lambda^2du}{\Lambda(x)^2},
\]
so $\Psi R_\varphi\le\tfrac12$ is exactly
\begin{equation}
  \int_0^x\!\Lambda(u)^2du \;\le\; \tfrac12\tanh\tfrac{x}{2}\,\Lambda(x)^2,
  \tag{D.8}
\end{equation}
and the same reduction turns $\Psi R_{GE}\le1$ into
\begin{equation}
  \int_0^x\! u\,\Lambda(u)\,du \;\le\; x\tanh\tfrac{x}{2}\,\Lambda(x) .
  \tag{D.9}
\end{equation}
The half-angle tangent is not an accident: $\tanh(x/2) =
(\cosh x - 1)/\sinh x = \Lambda'(x)/\sinh x$, so both statements
compare an integral of $\Lambda$ against $\Lambda$ and its own
derivative.

\textbf{Where $\tfrac12$ and $1$ come from.} Both are the $x\to\infty$
limits, and each is one application of l'H\^opital's rule. Written as
above, both products have the form $\coth(x/2)\int_0^x\!f/f(x)$ --- with
$f = \Lambda^2$ for the gravity channel and $f = u\Lambda$ for the
ground-effect channel --- and $\coth(x/2)\to1$. In the remaining ratio
numerator and denominator both diverge, so
\[
  \lim_{x\to\infty}\frac{\int_0^x\!f(u)\,du}{f(x)}
  = \lim_{x\to\infty}\frac{f(x)}{f'(x)}
  = \Bigl(\lim_{x\to\infty}\frac{f'}{f}\Bigr)^{\!-1},
\]
and the two logarithmic derivatives are elementary:
\[
  \frac{f'}{f}\bigg|_{f=\Lambda^2} = \frac{2\Lambda'}{\Lambda}
  = \frac{2(\cosh x - 1)}{\sinh x - x}\;\longrightarrow\;2,
  \qquad
  \frac{f'}{f}\bigg|_{f=u\Lambda} = \frac1u + \frac{\Lambda'}{\Lambda}
  \;\longrightarrow\;1,
\]
giving $\tfrac12$ and $1$. Each constant is thus the reciprocal of its
kernel's exponential growth rate, which is what $\int_0^xe^{ku}du =
e^{kx}/k$ says: the gravity channel gets the smaller constant because
its kernel, $\Lambda^2\sim e^{2x}/4$, is the square of the other's,
$u\Lambda\sim u\,e^{x}/2$.

\textbf{The only theorem used.} Everything below rests on one
corollary of the mean value theorem, applied four times in all --- three
for (D.8) and once for (D.9) --- and nowhere supplemented.

\emph{Lemma (a corollary of the mean value theorem).} If $g$ is
continuous on $[0,b]$, differentiable on $(0,b)$ with $g'\ge0$ there,
and $g(0) = 0$, then $g\ge0$ on $[0,b]$. \emph{Proof:} for
$0\le s<t\le b$ the mean value theorem gives $g(t)-g(s) =
g'(\xi)(t-s)\ge0$, so $g$ is non-decreasing; taking $s=0$ gives
$g(t)\ge g(0)=0$. $\square$ (The symbol is $g$, not the $f$ of
Sec.~3: this lemma is applied to $\phi'$, $\phi$ and $D$, never to a
channel kernel.)

Note what is \emph{not} used. No limit as $x\to\infty$ enters: the
constant $\tfrac12$ comes from the algebra of (D.8), and the
inequality is established at each $x$ separately. No monotonicity of
$\Psi R_\varphi$ is assumed or proved. The limits quoted above serve
only to show the constants cannot be improved.

\textbf{Proof of (D.8).} Let $D(x)$ be the right side minus the left,
so that (D.8) is $D\ge0$. Every quantity divided by below ---
$\coth(x/2)$, $\Lambda^2$, $\Lambda$, $\cosh^2(x/2)$ --- is strictly
positive for $x>0$, so each step is an equivalence.

$D(0) = 0$, since $\tanh0 = \Lambda(0) = \int_0^0 = 0$. Differentiating
and factoring out $\Lambda>0$,
\[
  D'(x) = \Lambda(x)\,B(x),
  \qquad
  B(x) = \tfrac14\operatorname{sech}^2\tfrac{x}{2}\,\Lambda
         + \tanh\tfrac{x}{2}\,\Lambda' - \Lambda ,
\]
so $D'\ge0$ iff $B\ge0$. Write $y = x/2$, $S = \sinh y$, $C = \cosh y$,
giving $\Lambda = 2SC-2y$ and $\Lambda' = 2S^2$. Multiplying $B$ by
$2C^2>0$,
\[
  2C^2B = (SC-y) + 4S^3C - 4SC^3 + 4yC^2 ,
\]
and since $S^3C - SC^3 = SC(S^2-C^2) = -SC$ the two cubic terms
collapse to $-4SC$, leaving $y(4C^2-1) - 3SC = y(3+4S^2) - 3SC$. Now
$4S^2 = 2\cosh x - 2$ and $2SC = \sinh x$ and $y = x/2$, so
\[
  B(x) = \frac{\phi(x)}{4\cosh^2(x/2)},
  \qquad
  \phi(x) \triangleq x + 2x\cosh x - 3\sinh x ,
\]
and $B\ge0$ iff $\phi\ge0$. Finally
\[
  \phi'(x) = 1 - \cosh x + 2x\sinh x,
  \qquad
  \phi''(x) = \sinh x + 2x\cosh x .
\]

\emph{First application.} $g = \phi'$: both terms of
$\phi'' = \sinh x + 2x\cosh x$ are non-negative for $x\ge0$, and
$\phi'(0) = 1-1+0 = 0$. The Lemma gives $\phi'\ge0$.

\emph{Second application.} $g = \phi$: $\phi'\ge0$ by the first
application and $\phi(0) = 0$. The Lemma gives $\phi\ge0$, hence
$B\ge0$, hence $D' = \Lambda B\ge0$.

\emph{Third application.} $g = D$: $D'\ge0$ by the second application
and $D(0) = 0$. The Lemma gives $D\ge0$, which is (D.8). $\square$

\textbf{Proof of (D.9).} The same Lemma settles the ground-effect
channel, and here it is needed only once. Let $E$ be the right side of
(D.9) minus the left,
\[
  E(x) \;=\; x\tanh\tfrac{x}{2}\,\Lambda(x)
             \;-\; \int_0^x\! u\,\Lambda(u)\,du ,
  \qquad E(0) = 0 .
\]
Differentiating --- the integral contributes $x\Lambda(x)$ --- and
using $\operatorname{sech}^2(x/2) = 2/(\cosh x + 1)$ and
$\tanh(x/2) = \sinh x/(\cosh x + 1)$ to put everything over
$\cosh x + 1$, the hyperbolic terms collapse to
\begin{equation}
  E'(x) \;=\; \frac{\sinh^2x - 2x\sinh x + x^2\cosh x}{\cosh x + 1}
        \;=\; \frac{\Lambda(x)^2 + x^2\Lambda'(x)}{\cosh x + 1} ,
  \tag{D.10}
\end{equation}
the second equality being the identity $(\sinh x - x)^2 + x^2(\cosh x
- 1) = \sinh^2x - 2x\sinh x + x^2\cosh x$. The numerator is a sum of
two non-negative terms and the denominator is positive, so $E'\ge0$;
the Lemma gives $E\ge0$, which is (D.9). $\square$

No series is needed for either, and (D.9) costs one application of the
Lemma against three for (D.8) --- the completed square in (D.10) does
the work that $\phi$ did there.

No monotonicity of the products is claimed or needed --- only that they
do not exceed their limits, which is what (D.8) and (D.9) say.
Expanding $\rbar$ by (95),
\[
  \Psi\rbar = \bigl(\Psi R_\varphi\bigr)\rho_{\varphi,\max}
            + \bigl(\Psi R_{GE}\bigr)\rho_{GE,\max}
            \;\le\; \tfrac12\rho_{\varphi,\max} + \rho_{GE,\max},
\]
which is the left half of (97). \emph{$x$ has disappeared} --- and with
it $C_2$, $J_P$ and $\Wz$. This is the whole point of the exercise: the
bound no longer depends on the window length, so it does not have to be
solved for, only bounded well enough to keep $\rho_{GE,\max}$ finite
(which is what (92)--(93) do).

\section*{6. From the rate deviation to the threshold: second half of (97)}

The reported quantity is not $e_\omega$ but the onset moment, so the
deviation must be pushed one step further.

\textbf{How the onset responds to a perturbation.} The fit is
two-segment and \emph{both} segments are in the cost --- a constant
baseline before the onset, the closed-form rise after it, sharing that
baseline:
\begin{align*}
  J(t_c) &= \int_{t_a}^{t_c}\!\bigl(y - C_0\bigr)^2dt \\[2pt]
   &\quad+ \int_{t_c}^{t_{\rm end}}\!
     \Bigl(y - C_0 - C_1\bigl(\cosh C_2(t-t_c) - 1\bigr)\Bigr)^2dt ,
\end{align*}
with $C_0$ the pre-onset baseline, refitted at each candidate $t_c$
from the samples to its left. In the shifted variable $\tau = t - t_c$
the post-onset window is $[0,T]$ with $T = t_{\rm end}-t_c = \te$, the
window Sec.~2 averages over, so the $\rbar$ below is the same $\rbar$.

Write the data as the true pre-onset level, the nominal response from
the true onset, and a deviation,
\[
  y(t) = b + \hat\omega(t;t_c^*) + e_\omega(t),
  \qquad
  t_c = t_c^* + \delta t_c ,
\]
the perturbation being $\rho$ itself and not its average. Carrying $b$
rather than setting it to zero is what later gives $C_0$ something to
estimate; without it $C_0$ would enter the cost and leave the
stationarity condition unaccounted for.

One point about that definition is worth making explicitly, because
getting it wrong invalidates everything after it. \emph{$e_\omega$ is
anchored at the true onset $t_c^*$, not at the candidate.} It is the
difference between the recorded rate and the nominal solution of the
true problem, so once the run is recorded it is a fixed function of $t$
and
\[
  \frac{\partial e_\omega}{\partial t_c} \;=\; 0
  \qquad\text{identically.}
\]
When the residual is written $e_\omega - \delta t_c\,\chi - C$ below,
every appearance of the optimisation variable is in $\delta t_c$ and in
$C$; none is in $e_\omega$. Reading $e_\omega$ instead as ``the residual
at the current candidate'' would make it depend on $t_c$ and would
count the same dependence twice.

Name the sensitivity
\[
  \chi(t) \;\triangleq\;
  \frac{\partial\hat\omega(t;t_c)}{\partial t_c}\bigg|_{t_c = t_c^*}
  \;=\; -\,C_1C_2\sinh\bigl(C_2(t-t_c^*)\bigr) .
\]
It needs a symbol of its own, and not $\varphi$: that is the tilt angle
everywhere else, and $\langle e_\omega,\varphi\rangle$ would read as a
rate correlated against an angle. To first order
$\hat\omega(t;t_c^*+\delta t_c)\simeq\hat\omega(t;t_c^*) + \delta t_c\,\chi$,
so the post-onset residual is $e_\omega - \delta t_c\,\chi - (C_0-b)$:
the baseline enters only through its \emph{error}, a point taken up
below. Setting that error aside for the moment,
\[
  \frac{\partial J}{\partial t_c}
  = \underbrace{\bigl(y(t_c) - C_0\bigr)^2
              - \bigl(y(t_c) - C_0\bigr)^2}_{=\;0,\ \text{the two limits}}
    \;-\;2\!\int_{t_c}^{t_{\rm end}}\!\bigl(e_\omega - \delta t_c\,\chi\bigr)
    \chi\,dt
    \;-\;2\frac{dC_0}{dt_c}\!\int\! r \;=\; 0 .
\]
$t_c$ is the upper limit of one integral and the lower limit of the
other, so it produces two boundary terms, and they cancel
\emph{identically}: at $\tau = 0$ the rise contributes
$C_1(\cosh 0 - 1) = 0$, so the two-segment model is continuous at the
onset by construction and both terms are $\bigl(y(t_c)-C_0\bigr)^2$
with opposite signs. Nothing is discarded and no order argument is
needed. This is worth stating, because the first thing to check about a
fit whose window endpoint is also its optimisation variable is whether
the Leibniz terms were quietly dropped.

The baseline still couples the two segments, since $C_0$ depends on
$t_c$ through the samples it is fitted to. Its contribution to the
pre-onset integral is $(dC_0/dt_c)\int_{t_a}^{t_c}(y - C_0)$, which
vanishes --- for the least-squares constant by its normal equation, and
for the median the code actually uses because a median is piecewise
constant in the sample set, so $dC_0/dt_c = 0$ almost everywhere. What
could survive is the same drift entering the post-onset residual, but
$dC_0/dt_c = \bigl(y(t_c) - C_0\bigr)/(t_c - t_a)$ for the mean, and at
the onset the rate is still the baseline, so it vanishes there too.
Measured across the $140$ runs, $C_0$ moves by under
$7.4\times10^{-4}$~rad/s when the candidate onset shifts by one sample.
Dropping that term and solving,
\begin{equation}
  \delta t_c = \frac{\langle e_\omega,\chi\rangle}{\|\chi\|^2},
  \qquad
  \langle a,b\rangle \triangleq \int_{t_c^*}^{t_{\rm end}}\!a(t)\,b(t)\,dt ,
  \tag{D.11}
\end{equation}
the division legitimate because $\|\chi\|^2>0$ --- $\sinh(C_2\tau)$ is
not identically zero on a window of positive length.

\textbf{Where the second derivative of $\hat\omega$ went.} Only $\chi$
appears in the stationarity condition because $\partial J/\partial t_c$
differentiates the residual once; $\partial_{t_c}\chi$ belongs to the
\emph{second} derivative of $J$, and (D.11) is the Gauss--Newton form
that drops it. Writing the one-parameter Newton step exactly,
\[
  J'(t_c^*) = -2\langle e_\omega,\chi\rangle,
  \qquad
  J''(t_c^*) = 2\|\chi\|^2 - 2\langle e_\omega,\partial_{t_c}\chi\rangle,
\]
\[
  \delta t_c = -\frac{J'}{J''}
  = \frac{\langle e_\omega,\chi\rangle}
         {\|\chi\|^2 - \langle e_\omega,\partial_{t_c}\chi\rangle},
  \qquad
  \partial_{t_c}\chi = C_1C_2^{2}\cosh(C_2\tau) .
\]
The neglected term is the residual times a curvature, the one nonlinear
least squares always discards. It is legitimate here for a reason worth
stating plainly: \emph{the expansion parameter is $e_\omega$, not
$\delta t_c$}, and $\delta t_c$ is itself first order in $e_\omega$ by
(D.11). So $\langle e_\omega,\chi\rangle$ and
$\delta t_c\|\chi\|^2$ are both $\mathcal{O}(e_\omega)$ while
$\delta t_c\langle e_\omega,\partial_{t_c}\chi\rangle$ is
$\mathcal{O}(e_\omega^2)$ --- every term carrying a derivative of
$\chi$ picks up an extra factor of the deviation, including the
$(\delta t_c)^2\partial_{t_c}\chi$ left over from expanding the
residual itself.

\textbf{Solving for $\delta t_c$ directly.} Once the residual is
linearised, $J$ is an \emph{exact} quadratic in $\delta t_c$, so it is
cleaner to differentiate in that variable and solve for the value that
zeroes the derivative. Since $t_c = t_c^* + \delta t_c$ the operator is
the same, but the two approximations then sit in separate places:
linearising the residual is the only one, and the normal equation that
follows is exact for the quadratic it is applied to.

\textbf{What the baseline has to satisfy, and what it does not.} Keeping
$C_0$ inside that square matters. Differentiating
$(e_\omega - \delta t_c\chi - C_0)^2$ and solving gives
\begin{equation}
  \delta t_c = \frac{\langle e_\omega - C_0,\;\chi\rangle}{\|\chi\|^2}
  = \frac{\langle e_\omega,\chi\rangle}{\|\chi\|^2}
    \;-\; C_0\,\frac{\langle 1,\chi\rangle}{\|\chi\|^2},
  \tag{D.11a}
\end{equation}
and the second term does not vanish: $\chi$ is of one sign, so
$\langle1,\chi\rangle = -C_1(\cosh x - 1)\ne0$.

It does not follow that $C_0$ must be driven to zero, which is the
natural but wrong reading of the onset condition. Only the \emph{error}
of $C_0$ against the true pre-onset level $b$ survives: a run whose gyro
sits at $b = 0.2$~rad/s is displaced by nothing at all provided
$C_0 = b$, while a mismatch of $0.001$~rad/s on a run with $b = 0$
displaces the onset measurably. What is required is an unbiased estimate
of the pre-onset level --- which is why the estimator fits $C_0$ rather
than fixing it, and fits a median: the pre-onset gyro noise is
heavy-tailed, and a candidate that overshoots leaks the rise into the
segment being averaged.

Carrying (D.11a) through $\Delta M_{\rm crit} = \dot M\,\delta t_c$ and
$\|\chi\|^2 = C_1^2C_2\bigl[\tfrac14\sinh2x - \tfrac12 x\bigr]$, the
amplitude cancels as everywhere else and
\begin{equation}
  \Delta M_{{\rm crit},C}
  = \Wz\,(C_0 - b)\,
    \frac{\cosh x - 1}{C_2\bigl[\tfrac14\sinh 2x - \tfrac12 x\bigr]} .
  \tag{D.11b}
\end{equation}
Minimising the cost directly reproduces (D.11b) to three figures. Over
the box the gain runs from $0.042$~N\,m per rad/s at the slowest ramp to
$0.351$ at the fastest, the fast end being worse because the window is
shorter. The baselines the sweep actually forms --- medians over the
excitation window up to the candidate, not over the whole quiet record
before it --- are $3.5$~mrad/s in the median and $47.8$ at worst across
the $140$ runs, so bounding $|C_0-b|$ by $|C_0|$ gives $0.054$~mm of
offset in the median and $0.558$~mm at worst.

That bound is loose, and the direct measurement says so. Rerunning the
whole pipeline with $C_0$ pinned to zero instead of fitted, which is the
largest perturbation this term admits, moves the direction-paired
$M_{ff}$ by $0.80$~mN\,m in the median and $4.91$ at worst ---
$0.027$~mm and $0.163$~mm. The pairing cancels most of what the bound
forwards. Neither statistic is better: the ramp-invariance CV is
$0.0287$ against $0.0292$, a paired $t$ of $0.87$ on $19$ degrees of
freedom, and $66$ of the $140$ runs return an identical onset index
either way. The median is kept because a pinned zero has no defence
against genuine pre-onset motion, not because it measures better.

Numerically, $\langle e_\omega,\partial_{t_c}\chi\rangle/\|\chi\|^2$ is
$9.05\times10^{-5}$, $9.05\times10^{-4}$, $9.05\times10^{-3}$ and
$9.05\times10^{-2}$ as the deviation grows through four decades: exactly
proportional to it, as an $\mathcal{O}(e_\omega)$ relative error must
be. That ratio is the whole of the discrepancy reported below. With
$\Delta M_{\rm crit} = \dot M\,\delta t_c$ and $C_1 = \dot M/\Wz$ the
amplitude $C_1$ cancels between numerator and denominator, leaving
\begin{equation}
  \Delta M_{\rm crit}
  = -\,\frac{\Wz\,\langle e_\omega,\;\sinh(C_2\tau)\rangle}
            {C_2\,\|\sinh(C_2\tau)\|^2}.
  \tag{D.12}
\end{equation}
This is exactly the reduction \texttt{error\_budget.py} implements, and
it is the manuscript's (98)--(106) written with $\chi$ named:
(D.11) is its (104), the baseline term below is its (105)--(106).
Minimising $J$ numerically against a deviation of the true shape, over
four decades of deviation size, gives ratios to the formula of
$0.99995$, $0.99953$, $0.99530$ and $0.95559$: it is the first-order
term, as claimed. Minimising the full two-segment cost and the
post-onset term alone gives the same answer to five figures at every
size, which is the numerical form of the two paragraphs above.

\textbf{From (D.12) to (97b), step by step.} The reduction is short but
has three distinct moves in it, so they are separated here. Throughout,
$\tau$ runs from $0$ at the onset to $\te$ at the window end, and
$\langle a,b\rangle = \int_0^{\te}ab\,d\tau$.

\emph{Step 1: write the inner product out.} (D.12) reads
\[
  \Delta M_{{\rm crit},e}
  = -\,\frac{\Wz}{C_2}\,
    \frac{\int_0^{\te} e_\omega(\tau)\sinh(C_2\tau)\,d\tau}
         {\|\sinh(C_2\tau)\|^2}.
\]

\emph{Step 2: substitute the deviation.} $e_\omega$ is not free: (D.2)
gives it as the Duhamel integral of the forcing, with
$\dot h_\varphi(u) = \cosh(C_2u)/J_P$ from (87). Putting that in makes
the numerator a double integral over the triangle
$0\le s\le\tau\le\te$,
\[
  \int_0^{\te}\!\!\sinh(C_2\tau)
  \left[\int_0^{\tau}\!\frac{\cosh\bigl(C_2(\tau-s)\bigr)}{J_P}\,
        \rho(s)\,ds\right]dt .
\]

\emph{Step 3: exchange the order.} Fubini applies --- both integrands
are continuous on a bounded domain --- and the same triangle read the
other way has $s$ outermost with $\tau$ running from $s$ up to $\te$:
\[
  = \int_0^{\te}\!\rho(s)
    \underbrace{\left[\frac{1}{J_P}\int_s^{\te}\!\sinh(C_2\tau)
      \cosh\bigl(C_2(\tau-s)\bigr)\,d\tau\right]}
     _{\text{a function of }s\text{ alone}}ds .
\]
The forcing is now outside, multiplied by something that no longer
depends on it. That is the whole point of the step: what began as a
correlation of $e_\omega$ against a kernel has become a weighted
average of $\rho$ itself,
\begin{equation}
  \Delta M_{{\rm crit},e} = -\!\int_0^{\te}\!\rho(s)\,w(s)\,ds,
  \quad
  w(s) = \frac{\Wz\displaystyle\int_s^{\te}\!\sinh(C_2\tau)
                \cosh\bigl(C_2(\tau-s)\bigr)d\tau}
              {J_PC_2\|\sinh(C_2\tau)\|^2}.
  \tag{D.13}
\end{equation}

\emph{Step 4: three properties of the weight.} All three are checked in
code (\texttt{error\_budget.py -{}-check}).

\begin{enumerate}
\item[(i)] $w\ge0$: the integrand is a product of two hyperbolic
  functions of non-negative arguments.
\item[(ii)] $\int_0^{\te} w = 1$, \emph{exactly}. Integrate the
  numerator of (D.13) over $s$ and exchange the order once more --- the
  same triangle again ---
  \[
    \int_0^{\te}\!\!\int_s^{\te}\!\sinh(C_2\tau)
      \cosh\bigl(C_2(\tau-s)\bigr)d\tau\,ds
    = \int_0^{\te}\!\sinh(C_2\tau)
      \left[\int_0^{\tau}\!\cosh\bigl(C_2(\tau-s)\bigr)ds\right]d\tau .
  \]
  The inner integral is elementary, $u = \tau-s$ turning it into
  $\int_0^{\tau}\cosh(C_2u)du = \sinh(C_2\tau)/C_2$, so the whole
  numerator becomes $\|\sinh(C_2\tau)\|^2/C_2$ --- the very norm sitting
  in the denominator of (D.13). Hence
  \[
    \int_0^{\te} w
    = \frac{\Wz\,\|\sinh\|^2/C_2}{J_PC_2\|\sinh\|^2}
    = \frac{\Wz}{J_PC_2^{2}} = 1
  \]
  by $J_PC_2^{2} = \Wz$, the identity that removes the inertia
  everywhere else in this document. Neither the window length nor any
  constant survives. \texttt{error\_budget.py -{}-check} evaluates
  (D.13) at $\rho\equiv1$ and returns $-1.0005$ to $-1.0001$ across five
  operating points, which is quadrature error on an identity rather than
  evidence for it.
\item[(iii)] $w$ is non-increasing in $s$: as $s$ grows both the
  integration range and the $\cosh$ factor shrink.
\end{enumerate}

\emph{Step 5: what (i) and (ii) alone give.} A non-negative weight of
unit mass can do no more than pick out the largest value of what it
averages, so
\[
  \bigl|\Delta M_{{\rm crit},e}\bigr|
  \;\le\; \rmax\!\int_0^{\te}\! w \;=\; \rmax .
\]
This is already a result --- a uniform unmodelled moment displaces the
identified threshold one-for-one, as a consequence rather than an
assumption --- but it is the weak form, and over the box it is
$2.70$~mm of offset, larger than the accuracy actually achieved.

\emph{Step 6: what (iii) adds.} The forcing rises across the window
while the weight falls, so the two are oppositely monotone and
Chebyshev's inequality (D.4) applies a second time, in the same
direction as the first:
\[
  \Bigl|\int_0^{\te}\!\rho\,w\Bigr|
  \;\le\; \frac{1}{\te}\int_0^{\te}\!\rho \;\cdot\!\int_0^{\te}\! w
  \;=\; \rbar\cdot1 \;=\; \rbar .
\]
The maximum has been replaced by the window average, and by (ii) no
kernel gain is left over to multiply it.

\emph{Step 7: bound the average.} $\rbar$ was computed at (95) as
$R_\varphi\rho_{\varphi,\max} + R_{GE}\rho_{GE,\max}$, and the two
reduction factors are bounded there by $\tfrac17$ and $\tfrac15$.
Chaining Steps 6 and 7,
\begin{equation}
  \bigl|\Delta M_{{\rm crit},e}\bigr| \;\le\; \rbar
  \;\le\; \frac{\rho_{\varphi,\max}}{7} + \frac{\rho_{GE,\max}}{5},
  \tag{97b}
\end{equation}
the right half of (97). Note the direction of the second inequality:
$\rbar$ is the average, and the right-hand side bounds it --- it is not
a definition of $\rbar$.

\emph{The same bound in two lines, without Fubini.} Steps 1--6 build the
weight in order to expose its three properties, and (iii) in particular
is what licenses Chebyshev. But if only the \emph{value} of the bound is
wanted, there is a shorter route that uses nothing beyond (90) and the
projection (D.11), and it is worth recording because it answers the
obvious question --- whether the pointwise bound on $e_\omega$ can be
used directly --- with a yes.

Evaluate (90) at each $\tau$ rather than only at the window end. Since
$\rho$ rises, its running average does too, so $\rbar(\tau)\le\rbar$ and
\begin{equation}
  \bigl|e_\omega(\tau)\bigr| \;\le\; E(\tau)
   \;\triangleq\; \frac{\rbar\,\sinh C_2\tau}{J_PC_2}
  \qquad\text{for every }\tau\in[0,\te].
  \tag{D.13a}
\end{equation}
The onset shift is the projection (D.11),
$\delta t_{\rm c} = \langle e_\omega,\chi\rangle/\|\chi\|^{2}$ with
$\chi = -C_1C_2\sinh C_2\tau$, and $\chi$ is of one sign, so the
absolute value may be taken inside:
\begin{equation}
  \bigl|\delta t_{\rm c}\bigr|
  \;\le\; \frac{\langle E,|\chi|\rangle}{\|\chi\|^{2}} .
  \tag{D.13b}
\end{equation}
Now the point. \emph{$E$ and $\chi$ are both proportional to
$\sinh C_2\tau$}, so the integral $\int_0^{\te}\sinh^2 C_2\tau\,d\tau$
appears in numerator and denominator alike and cancels without ever
being evaluated:
\[
  \frac{\langle E,|\chi|\rangle}{\|\chi\|^{2}}
  = \frac{\bigl(\rbar C_1/J_P\bigr)\int_0^{\te}\sinh^2 C_2\tau\,d\tau}
         {C_1^{2}C_2^{2}\int_0^{\te}\sinh^2 C_2\tau\,d\tau}
  = \frac{\rbar}{J_PC_1C_2^{2}}
  = \frac{\rbar}{C_1\Wz},
\]
using $J_PC_2^{2} = \Wz$ at the last step. Carrying it through
$\Delta M_{{\rm crit},e} = \dot M\,\delta t_{\rm c}$ and
$C_1 = \dot M/\Wz$, so that $C_1\Wz = \dot M$,
\begin{equation}
  \bigl|\Delta M_{{\rm crit},e}\bigr|
   \;\le\; \dot M\cdot\frac{\rbar}{C_1\Wz}
   \;=\; \rbar ,
  \tag{D.13c}
\end{equation}
which is Step 6 again. Checked numerically over a range of
$(C_2,J_P,C_1,\te)$ the ratio of the two is $1$ to six decimals.

The two routes are the same inequality seen twice: what appears as
$\int w = 1$ in the long route appears as the cancellation of
$\int\sinh^2$ in the short one. The long route is still the one to
keep, because it exhibits $w$ and its monotonicity, and without those
Chebyshev has nothing to act on --- Step 6 would have no justification
and the bound would stop at $\rmax$, the $2.70$~mm form. The short route
is the check, and the answer to a reader who asks why the deviation
bound cannot simply be inserted into the budget: inserted through the
projection, it can, and it gives exactly $\rbar$.

\emph{The band, and what the minimiser can reach.} (D.13a) is a
statement about the \emph{data}: the recorded rate, less its baseline,
lies within $\pm E(\tau)$ of the nominal cosh. The natural question is
whether the same is true of the \emph{answer} --- the minimiser returns
a curve, and one wants to know whether that curve stays in the band the
data lies in.

It does, but the reader should be told at once in what sense, because
the derivation below is a \emph{corollary of} (D.13c) and not an
independent check on it. It takes $|\delta| \le \rbar/\dot M$ from
(D.13c) and multiplies by $C_1C_2\sinh$; $E$ was itself built from the
same $\rbar$, and $\Wz = J_PC_2^{2}$ makes the two agree exactly. The
calculation could not have come out any other way, and nothing would
have been learned had it come out differently. It is worth writing
because it converts a statement about a number into a statement about a
curve --- but the load-bearing question, whether a fitted curve must lie
near the truth \emph{at all}, is not answered by it. That question is
general, has a short answer, and is taken up after this.

With the amplitude pinned at $C_1 = \dot M/\Wz$, the fitted curve is the
nominal one with its onset moved and nothing else:
\begin{equation}
  \hat\omega(\tau) = C_1\bigl(\cosh C_2(\tau - \delta) - 1\bigr),
  \qquad \delta \triangleq \delta t_{\rm c} .
  \tag{D.13d}
\end{equation}
Subtracting the nominal and applying
$\cosh A - \cosh B = 2\sinh\frac{A+B}{2}\sinh\frac{A-B}{2}$ --- an
identity, no expansion yet ---
\begin{equation}
  \hat\omega(\tau) - \omega_{\rm nom}(\tau)
   = -2C_1\sinh\frac{C_2\delta}{2}\,
       \sinh C_2\Bigl(\tau - \frac{\delta}{2}\Bigr).
  \tag{D.13e}
\end{equation}
Expanding the second factor and writing $g \triangleq C_2|\delta|$,
so that $2\sinh\frac g2\cosh\frac g2 = \sinh g$ and
$2\sinh^2\frac g2 = \cosh g - 1$,
\begin{equation}
  \bigl|\hat\omega - \omega_{\rm nom}\bigr|
   \;\le\; C_1\sinh g\,\sinh C_2\tau
     \;+\; C_1\bigl(\cosh g - 1\bigr)\cosh C_2\tau .
  \tag{D.13f}
\end{equation}
This is exact. To first order in $g$ the second term drops and
$\sinh g \to g$, and then the chain closes on itself:
\begin{equation}
\begin{aligned}
  \bigl|\hat\omega(\tau) - \omega_{\rm nom}(\tau)\bigr|
  &\;\le\; C_1C_2|\delta|\,\sinh C_2\tau
   \;\le\; \frac{\rbar}{C_1\Wz}\,C_1C_2\sinh C_2\tau \\[2pt]
  &\;=\; \frac{\rbar\,C_2}{\Wz}\sinh C_2\tau
   \;=\; \frac{\rbar\,\sinh C_2\tau}{J_PC_2}
   \;=\; E(\tau),
\end{aligned}
  \tag{D.13g}
\end{equation}
the middle inequality being (D.13b)--(D.13c) and the last step
$\Wz = J_PC_2^{2}$ again. \textbf{The fitted cosh lies inside the very
band that contains the data}, and inside it with the same envelope
$E(\tau)$, not a larger one.

The inequality in (D.13g) is an equality when $|\delta|$ takes its
extreme admissible value $\rbar/\dot M$. So the band is not a container
that happens to hold the fit: to first order it is exactly the tube
swept by the admissible onset family
\[
  \Bigl\{\,C_1\bigl(\cosh C_2(\tau-\delta) - 1\bigr)
   \;:\; |\delta| \le \rbar/\dot M \,\Bigr\},
\]
whose two extreme members \emph{are} the band edges
$\omega_{\rm nom}\pm E$. Every curve the minimiser can return is in the
band, and every part of the band is reached by some such curve.

That makes the geometric statement and the threshold statement one and
the same. Since $\Delta M_{{\rm crit},e} = \dot M\delta$, the assertions
``$|\Delta M_{{\rm crit},e}| \le \rbar$'' and ``the fitted cosh lies
between $\omega_{\rm nom} \pm \rbar\sinh(C_2\tau)/(J_PC_2)$'' are two
readings of a single inequality, one about the answer and one about the
curve. Nothing beyond (D.13a) and the projection was used to pass
between them.

Two qualifications. First, the pinning is doing real work here. Were
$C_1$ free as well, the fit would be $R\cosh(C_2\tau+\psi) - C_1$ with
$R\neq C_1$, and the amplitude channel is governed by
$P = \int\cosh(C_2s)\rho\,ds$ rather than by $Q$; $P$ carries no factor
that vanishes at $s = 0$ and is not bounded by $\rbar$, so the fitted
curve could then leave the band. It is precisely because
$C_1 = K\dot M$ is fixed that the only freedom left is the one $E$
bounds. Second, (D.13g) is first order in $g$, and (D.13f) says what
that costs. Over the operating box, with $\rbar = 12.04$~mN\,m and
$C_2 = 5.05$:

\begin{center}\small
\begin{tabular}{@{}rrrrr@{}}
\toprule
$\dot M$ & $|\delta|_{\max} = \rbar/\dot M$ & $g = C_2|\delta|_{\max}$
 & (D.13f)$/E$ at $x=3$ & dropped term at $\tau=0$ \\
N\,m/s & s & & & \% of peak rate \\
\midrule
$0.10$ & $0.1204$ & $0.608$ & $1.377$ & $2.10$ \\
$0.20$ & $0.0602$ & $0.304$ & $1.169$ & $0.51$ \\
$0.45$ & $0.0268$ & $0.135$ & $1.071$ & $0.10$ \\
$0.80$ & $0.0150$ & $0.076$ & $1.039$ & $0.03$ \\
$1.20$ & $0.0100$ & $0.051$ & $1.026$ & $0.01$ \\
\bottomrule
\end{tabular}
\end{center}

The exact band is wider than $E$ by $2.6\%$ at the fast ramp and $38\%$
at the slowest, and the widening is a second-order effect of an
\emph{a priori} onset displacement that is itself a fifth of the window
at $\dot M = 0.10$; the displacements actually measured are two orders
smaller, so in the campaign $g \lesssim 10^{-3}$ and the two forms are
indistinguishable. The last column is the one place the first-order
form fails outright: at $\tau = 0$ the envelope $E$ vanishes while
$C_1(\cosh g - 1)$ does not, so a curve whose onset moved early is
strictly outside the band there. What it is outside by is at most
$2.1\%$ of the peak nominal rate, and it is at the one end of the
window where nothing is being measured. Every claim above is verified
in \texttt{analysis/absorption\_check.py}.

\emph{The general statement, and what it costs.} Strip away every
special feature --- the cosh, the pinned amplitude, the hyperbolic
identities --- and ask the bare question. Let $M$ be \emph{any} set of
candidate curves in $L^2$ of the window, let the record be
$y = f + e$ with $f \in M$, and let
$\hat g \in \arg\min_{g\in M}\|y-g\|$. Then
\begin{equation}
  \|\hat g - f\| \;\le\; \|\hat g - y\| + \|y - f\|
                 \;\le\; \|f - y\| + \|y - f\|
                 \;=\; 2\|e\| ,
  \tag{D.14a}
\end{equation}
the first step the triangle inequality and the second the definition of
$\hat g$ together with $f \in M$ --- if $f$ is admissible, the minimiser
cannot do worse than $f$ does. That is the whole proof. It uses no
linearisation, no convexity, no identifiability, no smoothness, and no
property of $M$ whatsoever. So the answer to whether the fitted curve
must lie in a band around the nominal is \emph{yes, in general and
cheaply}, provided two things hold: the nominal is in the class, and the
minimum is global. Both are worth naming as assumptions rather than
left implicit, and both are examined below.

Three things are given up in exchange for that generality, and they are
exactly the three the earlier derivation had.

\emph{The constant is $2$, and it is sharp.} Take $M = \{f, g\}$, two
points, and $e = \tfrac12(g-f) + \varepsilon$ chosen so that $y$ falls
just on $g$'s side. The minimiser returns $g$, and
$\|\hat g - f\| = \|g-f\| \to 2\|e\|$. Nothing about (D.14a) can be
improved without assuming something about $M$. Assuming a little
recovers the factor. If $M$ is an affine subspace $f + V$ then
$\hat g - f = P_V e$ exactly, and orthogonal projection onto a subspace
is non-expansive, so
\begin{equation}
  \|\hat g - f\| = \|P_V e\| \;\le\; \|e\| ,
  \tag{D.14b}
\end{equation}
with equality iff $e \in V$. For a curved $M$ of reach $r$ in the sense
of Federer --- the largest radius within which nearest-point projection
is single-valued --- the projection is Lipschitz with constant
$r/(r-d)$ at distance $d$, giving
\begin{equation}
  \|\hat g - f\| \;\le\; \frac{\|e\|}{1 - \|e\|/r}
  \qquad (\|e\| < r),
  \tag{D.14c}
\end{equation}
which is (D.14b) as $\|e\|/r\to0$ and (D.14a) at $\|e\| = r/2$. Our
class is curved, so the honest constant is (D.14c); measured on it,
$r$ runs from $0.068$ at $\dot M = 0.10$ to $3.35$ at $1.20$, against
$\|E\| = 0.0174$, so the constant is $1.348$ at the slowest ramp and
$1.005$ at the fastest. An adversarial search over $602$ profiles with
$|e|\le E$ --- including $e = \pm E$ and random sign-switching ones ---
finds a worst realised ratio of $1.0005$: just above the affine value,
as a curved class requires, and far inside both bounds.

\emph{It is an $L^2$ statement, not a pointwise one.} (D.14a) bounds
$\|\hat g - f\|$, not $|\hat g(\tau) - f(\tau)|$, and no amount of work
on it will produce the second, because a small $L^2$ error is
compatible with a large excursion on a short interval. The pointwise
band of (D.13g) is available only because the class is one-parameter
and (D.13e) says exactly what a member looks like. This is the real
division of labour: (D.14a) is what guarantees closeness, and the
parametrisation is what converts closeness into a band.

\emph{What the direction buys.} (D.14a) uses only the \emph{size} of
$e$. The projection route (D.13b) uses $\langle e,\chi\rangle$ --- where
$e$ points, not merely how large it is. Here that turns out to buy
nothing, and for an instructive reason: $E \propto \sinh C_2\tau$ and
$\chi \propto \sinh C_2\tau$ are parallel, so Cauchy--Schwarz is tight
and the two agree identically. Numerically,
\[
  \frac{\langle E,|\chi|\rangle}{\|\chi\|^{2}}
  \;=\; \frac{\|E\|}{\|\chi\|}
  \;=\; \frac{\rbar}{\dot M}
\]
to seven decimals at every ramp rate. So (D.13b), (D.14b) and (D.13c)
are three routes to one number, and the arrangement noted at the head of
this passage was not a trick --- but it was not doing the work either.

\emph{Where the argument actually fails.} Not in the constant, and not
in the geometry, but in $f \in M$. Nothing above survives if the class
does not contain the nominal; the general form is
\begin{equation}
  \|\hat g - f\| \;\le\; 2\|e\| + d_M ,
  \qquad d_M \triangleq \operatorname{dist}(f, M),
  \tag{D.14d}
\end{equation}
by the same two lines with $\|f-y\|$ replaced by the distance to the
nearest admissible curve. Since $C_2$ is pinned per configuration, a
misidentified $C_2$ is precisely this failure --- and it is not
hypothetical, the fitted values across the campaign spanning $3.50$ to
$8.00$. Fitting noiseless data generated at $C_2^{\rm true}$ with the
model pinned at $5.046$, at $\dot M = 1.20$~N\,m/s:

\begin{center}\small
\begin{tabular}{@{}rrrrrr@{}}
\toprule
$C_2^{\rm true}$ & error & $d_M$ & $d_M/2\|e\|$
 & onset bias & threshold bias \\
 & \% & & & ms & mN\,m \\
\midrule
$4.100$ & $-18.7$ & $0.0102$ & $0.38$ & $+36.3$ & $+43.5$ \\
$4.600$ & $-8.8$  & $0.0053$ & $0.18$ & $+18.0$ & $+21.6$ \\
$5.046$ & $0.0$   & $0.0000$ & $0.00$ & $0.0$   & $0.0$ \\
$5.600$ & $+11.0$ & $0.0083$ & $0.20$ & $-24.5$ & $-29.4$ \\
$6.400$ & $+26.8$ & $0.0246$ & $0.44$ & $-64.2$ & $-77.1$ \\
$8.000$ & $+58.5$ & $0.0837$ & $0.82$ & $-150.0$ & $-180.0$ \\
\bottomrule
\end{tabular}
\end{center}

A $C_2$ wrong by $9\%$ biases the identified threshold by $21.6$~mN\,m
with no noise present at all --- $1.8\times$ the entire a priori budget
$\rbar = 12.04$~mN\,m --- and at $58\%$ the fit saturates the $\pm150$~ms
search limit. This is a bias, not a bound: it does not average down over
runs, and no error model on $\rho$ covers it.

\emph{The other two directions.} A reader who counts parameters will
object that the model has four of them and this appendix has projected
onto two. The objection is right and the answer is that they belong to
different stages. Differentiating
$g = C_1\bigl(\cosh C_2(\tau-\delta)-1\bigr)+c$ at the nominal gives
four tangent directions,
\begin{equation}
\begin{aligned}
  u &= \partial g/\partial C_1 = \cosh C_2\tau - 1, &\quad
  \chi &= \partial g/\partial\delta = -C_1C_2\sinh C_2\tau, \\
  \mathbf{1} &= \partial g/\partial c, &\quad
  v &= \partial g/\partial C_2 = C_1\tau\sinh C_2\tau ,
\end{aligned}
  \tag{D.15a}
\end{equation}
and only the last of them leaves $\operatorname{span}\{\cosh,\sinh,1\}$
--- it carries the factor $\tau$, so it is the resonant direction of
\S4 and the one the shift identity cannot reach.

The onset sweep holds $C_1$ and $C_2$ fixed, so at that stage $u$ and
$v$ are not available and the projection onto
$\operatorname{span}\{\chi,\mathbf{1}\}$ is exactly right. What the two
frozen directions contribute is not absorption but \emph{bias}: if the
calibration supplies $C_1(1+\varepsilon_1)$ and $C_2(1+\varepsilon_2)$,
the deviation the sweep must accommodate is not $e_\omega$ but
$e_\omega - \varepsilon_1u - \varepsilon_2C_2v$, and the normal equation
(D.11) picks up two extra terms. Both integrate in closed form. With
$x = C_2\te$,
\begin{equation}
\begin{aligned}
  \langle u,\chi\rangle &= -C_1^{2}A(x)\big/1, &\quad
  \langle v,\chi\rangle &= -C_1^{2}D(x)/C_2, \\
  \|\chi\|^{2} &= C_1^{2}C_2B(x),
\end{aligned}
  \tag{D.15b}
\end{equation}
\begin{equation}
\begin{aligned}
  A(x) &= \tfrac14\cosh 2x - \cosh x + \tfrac34, \qquad
  B(x) = \tfrac14\sinh 2x - \tfrac12x, \\
  D(x) &= \tfrac14x\sinh 2x - \tfrac18\cosh 2x + \tfrac18
          - \tfrac14x^{2},
\end{aligned}
  \tag{D.15c}
\end{equation}
so that, carrying each through $\Delta M = \dot M\delta$ and
$C_1\Wz = \dot M$,
\begin{equation}
  \Delta M_{{\rm crit},\varepsilon}
   = \varepsilon_1\,\frac{\dot M}{C_2}\,\frac{A(x)}{B(x)}
   \;+\; \varepsilon_2\,\frac{\dot M}{C_2}\,\frac{D(x)}{B(x)} .
  \tag{D.15d}
\end{equation}
Both scale with $\dot M/C_2$, the moment swept in one plant time
constant, and $D/B$ exceeds $A/B$ by a factor $2.4$ to $3.6$ over the
operating windows: \emph{a misidentified exponent is about three times
as costly as a misidentified amplitude}. At $x = 3.735$ and
$\dot M = 0.45$~N\,m/s, (D.15d) gives $0.82$ and $2.90$~mN\,m per $1\%$
against $1.30$ and $3.72$ measured on the full nonlinear fit --- the
gap being the linearisation and the pre-onset segment, which the closed
form does not carry. $A$ and $D$ agree with direct quadrature to
machine precision.

Two things keep this from entering the reported budget. The
calibration is per configuration and carries no direction split, so
$\varepsilon_1$ and $\varepsilon_2$ are common to the two runs and
(D.15d) is antisymmetric in $\dot M$; measured on the full nonlinear
tip-over the half-sum rejects $99.77$ to $99.98\%$ of it, leaving
$0.000$~mm at $\pm10\%$ amplitude and $\pm5\%$ exponent error. And what
does survive goes to the \emph{weight} of (34), which reads the
difference rather than the sum, at $0.027\%$ of $W$ per $1\%$ of
amplitude error and $0.077\%$ per $1\%$ of exponent error. Verified in
\texttt{analysis/two\_sided\_cancellation.py}.

That is the useful conclusion of this passage, and it is not the one the
band statement suggests. The band argument is safe but nearly empty; the
sensitivity that matters sits in the assumption it takes for granted.
It also says what the per-configuration PNLS calibration is \emph{for}:
not to improve the fit cosmetically, but to buy $f \in M$, which is the
hypothesis every bound in this appendix rests on. Verified in
\texttt{analysis/general\_band.py}.

\begin{center}\small
\begin{tabular}{@{}rll@{}}
\toprule
step & what changes & why it is allowed \\
\midrule
1--2 & $e_\omega \to$ Duhamel double integral & (D.2), the exact solution \\
3 & double integral $\to \int\rho\,w$ & Fubini on the triangle \\
4 & --- & $w\ge0$, $\int w = 1$, $w\downarrow$ \\
5 & $\int\rho w \to \rmax$ & (i) and (ii) \\
6 & $\rmax \to \rbar$ & (iii) plus Chebyshev (D.4) \\
7 & $\rbar \to \rho_{\varphi}/7 + \rho_{GE}/5$ & (95) \\
\bottomrule
\end{tabular}
\end{center}


\textbf{Why the constants differ from the left half.} The left half
asks how far the rate curve is displaced, so the kernel integral
$\Psi$ multiplies $\rbar$ and the products $\Psi R$ climb to
$\tfrac12$ and $1$. The right half asks how far the \emph{minimiser}
moves, and the minimiser is found by correlating against
$\sinh(C_2\tau)$ --- a normalised operation, $\int w = 1$, with no
kernel gain left over. So there $\rbar$ enters bare, and the constants
are the $R$'s themselves, $\tfrac17$ and $\tfrac15$. The threshold
bound is the stronger of the two by roughly a factor of four, which is
why it, and not the relative rate bound, is the number carried into the
offset budget.

\section*{7. The threshold bound by hand}

The manuscript reports only the threshold bound, and that branch uses
neither $\Psi$ nor $x$: the weight of (D.13) is normalised, so $\rbar$
enters bare and the constants are the $R$'s themselves. What is left is
arithmetic. No hyperbolic function appears below, $C_2$ never enters,
and the window length is never solved for --- $\tfrac17$ and
$\tfrac15$ are suprema over \emph{all} $x$, so no $x$ has to be known
to use them.

Given the weight $W$, the height $z_{CoM}$, the arm
$a = l_p + p_{\rm off}$, the mass $m$, a CAD inertia $J_{CoM}$, the
ramp rate $\dot M$, the tilt cap $\varphi_{\max}$ and the ground-effect
slope $\beta_M$, six lines:
\begin{enumerate}
\item[(1)] $J_P \le J_{CoM} + m\bigl(z_{CoM}^2 + a^2\bigr)$
  \hfill --- (82) at its upper end
\item[(2)] $\Delta M_{\rm win} \le
  \bigl(6\varphi_{\max}J_P\dot M^{2}\bigr)^{1/3}$ \hfill --- (93)
\item[(3)] $\rho_{\varphi,\max} = \tfrac12 W a\,\varphi_{\max}^{2}$
  \hfill --- the gravity remainder at the cap
\item[(4)] $\rho_{GE,\max} =
  |\beta_M|\,\Delta M_{\rm win}\,\varphi_{\max}$
\item[(5)] $|\Delta M_{\rm crit}| \le
  \rho_{\varphi,\max}/7 + \rho_{GE,\max}/5$ \hfill --- (97b)
\item[(6)] divide by $W_{\min}$ for the equivalent offset.
\end{enumerate}

Over the geometric box of Sec.~VI-D, with $W = 31.59$~N,
$z_{CoM} = 0.30$~m, $m = 3.221$~kg, $J_{CoM} = 0.0537$~kg\,m$^2$,
$\varphi_{\max} = 10^\circ$, $\beta_M = 0.0345$~rad$^{-1}$ and
$W_{\min} = 30.08$~N:

\begin{center}\small
\begin{tabular}{@{}llrrrrrr@{}}
\toprule
axis & $\dot M$ & $J_P$ & $\Delta M_{\rm win}$
 & $\rho_{\varphi,\max}$ & $\rho_{GE,\max}$
 & $|\Delta M_{\rm crit}|$ & offset \\
 & N\,m/s & kg\,m$^2$ & N\,m & mN\,m & mN\,m & mN\,m & mm \\
\midrule
roll $(a = 0.160)$  & $0.10$ & $0.4260$ & $0.165$ & $76.98$ & $0.99$
  & $11.20$ & $0.372$ \\
                    & $0.45$ & $0.4260$ & $0.449$ & $76.98$ & $2.70$
  & $11.54$ & $0.384$ \\
                    & $1.20$ & $0.4260$ & $0.863$ & $76.98$ & $5.20$
  & $12.04$ & $0.400$ \\
\midrule
pitch $(a = 0.130)$ & $0.10$ & $0.3980$ & $0.161$ & $62.55$ & $0.97$
  & $\phantom{0}9.13$ & $0.304$ \\
                    & $0.45$ & $0.3980$ & $0.439$ & $62.55$ & $2.64$
  & $\phantom{0}9.46$ & $0.315$ \\
                    & $1.20$ & $0.3980$ & $0.844$ & $62.55$ & $5.08$
  & $\phantom{0}9.95$ & $0.331$ \\
\bottomrule
\end{tabular}
\end{center}

Three readings.

\emph{The inertia hardly matters.} $J_P$ enters only through
$\Delta M_{\rm win}$, which enters only through $\rho_{GE,\max}$, which
is $2$--$10\%$ of the total; and it enters as a cube root. Halving
$J_P$ moves the bound by $-1.8\%$, doubling it by $+2.2\%$. A crude
inertia is therefore good enough, which is worth knowing because (82)
is the one input here that is not directly measured.

\emph{Line (5) is where the work went.} Properties 1 and 2 of the
weight alone --- $w\ge0$ and $\int w = 1$ --- already give
$|\Delta M_{\rm crit}|\le\rho_{\max} = \rho_{\varphi,\max} +
\rho_{GE,\max}$, without Chebyshev and without any $R$. That is
$2.70$~mm at worst over the box, which \emph{exceeds} the $1.64$~mm
validation RMS: a guarantee weaker than the measured performance, and
so of no use. The factor of seven that $\tfrac17$ and $\tfrac15$ supply
is what brings it to $0.40$~mm and clears the measurement by $4.2\times$.

\emph{The rate bound is not needed for this.} The same construction
also bounds the relative rate deviation, with $\tfrac12$ and $1$ in
place of $\tfrac17$ and $\tfrac15$ and a factor $\Psi = x\coth(x/2)$;
that is Secs.~4--5, it is weaker by about a factor of four, and no
number above depends on it.

\section*{8. A worked example, checkable by hand}

\emph{On which $x$ to use.} Two values circulate below and they must not
be confused. The cubic shortcut (92) supplies $\bar x =
(6\Lambda_{\max})^{1/3}$, an \emph{upper bound} on $x$ obtained by
replacing $\Lambda(x)$ with its leading term $x^3/6$. The exact value
solves
\begin{equation}
  \Lambda(x) \;=\; \Lambda_{\max}
  \;\triangleq\; \frac{\varphi_{\max}\Wz C_2}{\dot M},
  \tag{D.14}
\end{equation}
which has a unique root because $\Lambda$ increases strictly from $0$.
The shortcut exists only to spare (97b) a transcendental solve it does
not need --- there $x$ appears nowhere, and $\Delta M_{\rm win}$ enters
through a channel worth under a tenth of the total. Wherever $x$ itself
appears, (D.14) is solved instead, for two reasons. It is loose: at the
slowest ramp $\Lambda$ exceeds $x^3/6$ by $3.6\times$, and anything
downstream of $\sinh$ exponentiates that. And it is not automatically
valid: substituting an upper bound into $R(x)\sinh x$ presumes that
product monotone, which is true but is not among the facts Sec.~5
proves, whereas solving (D.14) presumes nothing.

Both reduction factors close in finitely many terms, so the rows below
need no quadrature:
\begin{equation}
\begin{aligned}
  R_\varphi(x) &= \frac{A_\varphi(x)}{x\Lambda(x)^2}, &
  A_\varphi(x) &= \tfrac14\sinh 2x - \tfrac12 x
                  - 2x\cosh x + 2\sinh x + \tfrac13x^3, \\[2pt]
  R_{GE}(x) &= \frac{A_{GE}(x)}{x^{2}\Lambda(x)}, &
  A_{GE}(x) &= x\cosh x - \sinh x - \tfrac13x^3 ,
\end{aligned}
  \tag{D.15}
\end{equation}
each being $\int_0^x\Lambda^2$ and $\int_0^x u\Lambda$ integrated by
parts against $\int\sinh^2 = \tfrac14\sinh 2x - \tfrac12x$ and
$\int u\sinh u = u\cosh u - \sinh u$.

Take the roll axis at the fastest ramp, evaluated over the geometric
box of Sec.~VI-D rather than at any fitted constant: $W = 31.59$~N
(fully ballasted, $m = 3.221$~kg), $z_{CoM} = 0.30$~m, $a = l_p + p_{\rm off}
= 0.160$~m, $\dot M = 1.20$~N\,m/s, $\varphi_{\max} = 10^\circ =
0.17453$~rad, $\beta_M = 0.0345$~rad$^{-1}$, $J_{CoM}^{\rm CAD} =
0.0537$~kg\,m$^2$.

\begin{center}\small
\begin{tabular}{@{}llr@{}}
\toprule
step & expression & value \\
\midrule
geometry & $\Wz = Wz_{CoM}$ & $9.477$~N\,m\\
(82) lower & $J_P \ge m(z_{CoM}^2+a^2)$ & $0.3724$~kg\,m$^2$\\
(82) upper & $J_P \le J_{CoM}^{\rm CAD} + m(z_{CoM}^2+a^2)$ & $0.4261$~kg\,m$^2$\\
       & $C_2 \le \sqrt{g z_{CoM}/(z_{CoM}^2+a^2)}$ & $5.045$~s$^{-1}$\\
\midrule
tilt cap $\to\Lambda$ & $\Lambda_{\max} = \varphi_{\max}\Wz C_2/\dot M$ & $6.954$\\
(D.14) & solve $\sinh x - x = 6.954$ & $x = 2.9930$\\
       & \quad check: $\sinh(2.9930) - 2.9930$ & $6.9548$\\
(92)   & $\bar x = (6\Lambda_{\max})^{1/3}$, the shortcut & $3.4685$\\
(93)   & $\tau_{\rm end}\le(6\varphi_{\max}J_P/\dot M)^{1/3}$ & $0.719$~s\\
(93)   & $\Delta M_{\rm win} = \dot M\tau_{\rm end}$ & $0.863$~N\,m\\
\midrule
channel& $\rho_{\varphi,\max} = \tfrac12W a\,\varphi_{\max}^2$ & $76.98$~mN\,m\\
(93)   & $\rho_{GE,\max} = |\beta_M|\Delta M_{\rm win}\varphi_{\max}$ & $5.20$~mN\,m\\
\midrule
(D.15) & $\Lambda(x) = 6.9548$, $\Lambda(x)^2$ & $48.3690$\\
(D.15) & $A_\varphi(x)$ & $17.2171$\\
(D.6)  & $R_\varphi(x) = 17.2171/[(2.9930)(48.3690)]$ & $0.1189$\\
(D.15) & $A_{GE}(x)$ & $11.0389$\\
(D.7)  & $R_{GE}(x) = 11.0389/[(2.9930)^2(6.9548)]$ & $0.1772$\\
(95)   & $\rbar = R_\varphi\rho_{\varphi,\max}+R_{GE}\rho_{GE,\max}$
       & $10.08$~mN\,m\\
       & (against $\rmax = 82.2$~mN\,m) & \\
\midrule
(96)   & $\Psi(x) = x\coth(x/2)$ & $3.309$\\
(96)   & $\Psi\rbar/\Delta M_{\rm win}$ & $3.86\%$\\
(97)   & $(\tfrac12\rho_{\varphi,\max}+\rho_{GE,\max})/\Delta M_{\rm win}$
       & $5.06\%$\\
(97b)  & $\rho_{\varphi,\max}/7 + \rho_{GE,\max}/5$ & $12.04$~mN\,m\\
       & $\div\,W_{\min}$ & $0.400$~mm\\
\midrule
(D.3)  & envelope only, $\Psi\rmax/\Delta M_{\rm win}$ & $31.5\%$\\
\bottomrule
\end{tabular}
\end{center}

Read the rows in two groups. Everything from (D.15) down to (96) is
evaluated at the exact $x$ of (D.14), so each is the value the run
actually realises; the $\bar x$ row is carried only to show what the
shortcut costs, and using it would have given $\rbar = 9.54$ instead of
$10.08$~mN\,m --- an \emph{under}statement, since $R$ decreases. The
(97) and (97b) rows are of a different kind again: they contain no $x$
at all, so they hold whatever the window turns out to be. That is why
(97b), and not the $3.86\%$, is what the offset budget inherits.

Two readings. First, (97) costs about a third over (96) --- the price
of dropping $x$ from the statement --- whereas the envelope form costs
a factor of eight. Second, this is the \emph{worst admissible} case: the
box is evaluated at the $10^\circ$ tilt cap, at the top of the
$z_{CoM}$ range, and with the inertia at whichever end of (82) is
unfavourable for the quantity being bounded. Over the runs actually
flown the excursions are $5$--$9^\circ$ and the exact per-run
propagation, which convolves the measured $\rho$ instead of any
envelope, gives $0.04$~mN\,m in the median and $0.15$ at worst ---
$0.001$ to $0.004$~mm of offset.

\section*{9. Summary of the chain}

\begin{center}\small
\begin{tabular}{@{}rll@{}}
\toprule
& step & what it buys \\
\midrule
(D.2) & $e_\omega = \int\dot h_\varphi(\te-s)\rho(s)ds$
      & the exact deviation \\
(D.3) & $|\rho|\le\rmax$, integrate the kernel
      & (90); too weak by $\sim8\times$ \\
(D.4) & Chebyshev, $\rho\uparrow$ against $\dot h_\varphi(\te-\cdot)\downarrow$
      & $\rmax \to \rbar$ \\
(D.6--7) & known shapes $\Lambda^2$ and $u\Lambda$
      & the ratios $\rbar/\rmax$ \\
(D.6a) & $R = \int_0^1 f(xt)/f(x)dt$, $n_{\rm eff}\uparrow$
      & (95): $R_\varphi\le\tfrac17$, $R_{GE}\le\tfrac15$ \\
(96)  & divide by $\omega_{\rm nom,end}$, use $\Delta M_{\rm win}=\dot Mx/C_2$
      & $\Psi = x\coth(x/2)$; $\Wz$, $J_P$ gone \\
      & $\Psi R_\varphi\le\tfrac12$, $\Psi R_{GE}\le1$
      & (97) left half; $x$ gone \\
(D.12--13) & linearise the onset, Fubini, $w\ge0$, $\int w=1$
      & $\Delta M_{\rm crit} = -\langle\rho\rangle_w$ \\
      & Chebyshev again, $\rho\uparrow$ against $w\downarrow$
      & (97) right half: $\tfrac17,\tfrac15$ \\
\bottomrule
\end{tabular}
\end{center}

Chebyshev is used twice, for the same reason both times: the forcing is
concentrated at the end of the window and every kernel in the problem
weights the end of the window least. The two uses are independent ---
the first converts a kernel integral, the second a normalised
correlation --- which is why the constants they leave behind differ.

\end{document}
